\begin{abstract}
Los asistentes de domótica basados en modelos de lenguaje suelen depender de
servicios en la nube, lo que introduce alta latencia, costos recurrentes y
dependencia de conectividad. Los Small Language Models (SLMs) ejecutados
localmente constituyen una alternativa viable para hubs domóticos de bajo
costo, pero su capacidad para interpretar comandos en español rioplatense no
ha sido evaluada: el voseo, los diminutivos y las expresiones relativas de la
variedad (``bajale un toque'') no aparecen en los corpus de domótica en
inglés usados para entrenar y evaluar estos modelos, y su tratamiento no
puede suponerse transferido. Este trabajo evalúa doce SLMs abiertos de menos
de 2 mil millones de parámetros sobre un dataset de 32 comandos en español
rioplatense, con anotación de referencia de cinco campos (intención,
dispositivo, ubicación, valor y unidad). La verificación es automática en dos
etapas: comparación textual determinista campo a campo y un modelo de
lenguaje que actúa de juez, clasificando cada respuesta incorrecta contra una
taxonomía cerrada. Se reportan exactitud estricta y exactitud laxa
(semántica): su brecha cuantifica cuánto penaliza la coincidencia textual
exacta a respuestas que un evaluador humano consideraría correctas.

\keywords{modelos de lenguaje pequeños, domótica, evaluación automática,
LLM como juez, procesamiento del lenguaje natural, español rioplatense.}
\end{abstract}
